% 05_sensitivity.tex - which parameters the observations see; the first attempt
\section{What the observations can see}\label{sec:ident}
Before you calibrate, check that the observations depend on the parameters at all. Read this section. The models in it
were run for you; its one command only shows a file. The table changes one parameter at a time by 10\% and shows how the KGE of the streamflow, the head RMSE of
each well and the objective function change. The objective function is defined in \cmd{ostIn.txt}
(Step~\ref{st:ostin}): $-\mathrm{KGE}$ plus, for each well, its head RMSE times a weight $1/(2\sigma)$, where $\sigma$ is the
standard deviation of that well's observed heads. The weights are \V{t6:w101}, \V{t6:w102} and \V{t6:w103} for OW\_NEAR,
OW\_UPPER and OW\_SLOPE. With the true parameters the KGE is \V{t6:sens:kge}, the RMSEs are \V{t6:sens:rmse101},
\V{t6:sens:rmse102} and \V{t6:sens:rmse103}~mm (the noise of the observations), and the objective is \V{t6:sens:base}.
\begin{center}\small
\begin{tabular}{llrrrrr}\toprule
 & & & \multicolumn{3}{c}{change of the head RMSE [mm]} & \\\cmidrule(lr){4-6}
Parameter & change & KGE & OW\_NEAR & OW\_UPPER & OW\_SLOPE & objective\\\midrule
$K_h$ GRAVEL & $-$10\% & $-$0.006 & $-$1 & 0 & +1 & +0.002\\
$K_h$ GRAVEL & +10\% & $-$0.005 & +6 & +4 & +1 & +0.073\\[2pt]
$S_y$ GRAVEL & $-$10\% & 0.000 & 0 & 0 & 0 & +0.004\\
$S_y$ GRAVEL & +10\% & 0.000 & 0 & 0 & 0 & $-$0.001\\[2pt]
$K_h$ SAND & $-$10\% & $-$0.002 & $-$1 & 0 & +19 & +0.003\\
$K_h$ SAND & +10\% & $-$0.002 & +1 & +1 & +15 & +0.023\\[2pt]
$K_h$ TILL & $-$10\% & 0.000 & 0 & 0 & +1\,845 & +0.776\\
$K_h$ TILL & +10\% & 0.000 & +1 & 0 & +1\,598 & +0.675\\[2pt]
\cmd{PERC\_COEFF} & $-$10\% & $-$0.017 & +2 & +1 & +220 & +0.122\\
\cmd{PERC\_COEFF} & +10\% & $-$0.017 & 0 & 0 & +185 & +0.090\\
\bottomrule
\end{tabular}

\end{center}
The true values are: $K_h$ of GRAVEL \V{t6:trueK_GRAVEL}~m/d (valley, layer~1); $S_y$ of GRAVEL \V{t6:trueSY_GRAVEL};
$K_h$ of SAND \V{t6:trueK_SAND}~m/d (valley, layer~2); $K_h$ of TILL \V{t6:trueK_TILL}~m/d (slopes); \cmd{PERC\_COEFF} of
SUBSOIL \V{t6:truePERC_SUBSOIL}~1/d. The table shows:
\begin{itemize}
\item The valley wells see the gravel. OW\_NEAR and OW\_UPPER respond to its conductivity, more to a higher value than to
  a lower one, and with their large weights this is most of its change of the objective. Its $S_y$ hardly shows.
\item Only OW\_SLOPE sees the sand, which carries the groundwater of the slopes down the valley. The RMSE there grows by
  \V{t6:sens:sandslopemin} to \V{t6:sens:sandslopemax}~cm, about the size of the noise. The small weight of OW\_SLOPE
  (\V{t6:w103}, set by the large spread of its heads) cuts that to hundredths of the objective. The valley wells hardly see
  the sand, although OW\_NEAR is screened in the sand 200~m from the pumping well PW\_1: its heads follow the creek and the
  gravel.
\item The till matters very much. The water table under the slopes lies at a height set by the ratio of recharge to
  conductivity. A 10\% change of the till's $K_h$ raises the RMSE at OW\_SLOPE to \V{t6:sens:K_TILL:lo:slope} and
  \V{t6:sens:K_TILL:hi:slope}~m.
\item \cmd{PERC\_COEFF} sets how much of the subsoil's drainage becomes recharge; the rest flows to the creek as interflow.
  So it shows in the streamflow and, through the recharge, at OW\_SLOPE too.
\end{itemize}
The objective is not exactly at its minimum at the true values. The observations are noisy, so a nearby parameter set can
fit them a little better (a negative change in the table). Read the table for the size of the changes on both sides.

\subsection*{A first attempt}
A first calibration searched four parameters with \V{t6:fa:maxiter} DDS runs: $K_h$ and $S_y$ of the gravel, $K_h$ of the
till, and \cmd{PERC\_COEFF}. Its files are in \cmd{\$WS/model/finished/first\_attempt}. Its parameters were:
\term{t6_fa_params.txt}
\begin{center}\small
\begin{tabular}{lrrrr}\toprule
Parameter & initial & true & found & error\\\midrule
$K_h$ of GRAVEL [m/d] & \V{t6:fa:initK_GRAVEL} & \V{t6:trueK_GRAVEL} & \V{t6:fa:bestK_GRAVEL} & \V{t6:fa:errK_GRAVEL}\%\\
$S_y$ of GRAVEL & \V{t6:fa:initSY_GRAVEL} & \V{t6:trueSY_GRAVEL} & \V{t6:fa:bestSY_GRAVEL} & \V{t6:fa:errSY_GRAVEL}\%\\
$K_h$ of TILL [m/d] & \V{t6:fa:initK_TILL} & \V{t6:trueK_TILL} & \V{t6:fa:bestK_TILL} & \V{t6:fa:errK_TILL}\%\\
\cmd{PERC\_COEFF} of SUBSOIL [1/d] & \V{t6:fa:initPERC_SUBSOIL} & \V{t6:truePERC_SUBSOIL} & \V{t6:fa:bestPERC_SUBSOIL} & \V{t6:fa:errPERC_SUBSOIL}\%\\\bottomrule
\end{tabular}\end{center}
Its objective weighted every well alike (\V{t6:headw}~times its head RMSE). It reached \V{t6:fa:cost} at run
\V{t6:fa:itBest}; with the true parameters it is \V{t6:fa:costtrue}. Two things went wrong, and the table above explains
both:
\begin{itemize}
\item The gravel's parameters ended where the objective is flat. With these observations, no number of runs pins them
  down.
\item DDS found a till conductivity that fitted OW\_SLOPE early (it kept the value of run \V{t6:fa:ktfrom} to the end),
  but together with a wrong \cmd{PERC\_COEFF}: less recharge and a slightly lower conductivity give the same heads on the
  slope. From there, changing either parameter alone costs much more (see the table) than the streamflow gains. So the
  search stays in that narrow valley of the objective, even over \V{t6:fa:maxiter} runs.
\end{itemize}
Two changes make the calibration well posed:
\begin{enumerate}
\item Calibrate what the data can see, and not two parameters that the data see only together. The calibration of this
  workshop searches the conductivity of the till, which the data see clearly, and that of the sand, which they see weakly,
  to show what happens to such a parameter. The gravel and \cmd{PERC\_COEFF} keep the values of Workshop~2, as if they were
  known from elsewhere (a pumping test, or a calibration of the streamflow alone). In this synthetic experiment those are
  their true values, so the test is whether DDS recovers the two conductivities.
\item Weight each well by the spread of its own observations. In the calibration years the heads of OW\_SLOPE range over
  \V{t6:swingm103}~m, those of OW\_NEAR and OW\_UPPER over \V{t6:swing101} and \V{t6:swing102}~cm. Weighted alike, as in the
  first attempt, OW\_SLOPE dominates the objective so much that DDS keeps any set that fits it. The sand and the till also
  trade off there: a higher sand conductivity drains the slopes, and a lower till conductivity raises them again. Measured
  against its own spread, each well and the streamflow get a comparable say.
\end{enumerate}
